%!TEX root = ../template.tex

\typeout{NT FILE chapter3.tex}%
% ============================================================
%  Chapter 3 — System Description and Methodology
% ============================================================
\chapter{System Description and Methodology}
\label{chap:system-description}

This chapter presents the technical substrate on which the research contributions are built and the methodological approach that governs their design, implementation, and evaluation. Section~\ref{sec:system-description} describes the Closer GenAI Framework — referred to throughout this work as Maestro — in its pre-intervention state, characterising its architectural structure, data flows, and the specific gaps that motivate each research objective. Section~\ref{sec:methodology} formalises the research methodology, detailing how each contribution is designed, implemented, and evaluated against the success criteria stated in Section~\ref{sec:objectives}.

% ============================================================
\section{System Description}
\label{sec:system-description}
% ============================================================

\subsection{Overview and Architectural Pattern}
\label{subsec:overview}

Maestro is a multi-repository, Azure-native, microservices-based Retrieval-Augmented Generation platform designed to support the deployment of enterprise chatbots, and related ~\cite{closer2025maestro}. Its architectural pattern is a REST/JSON service mesh: independent FastAPI services communicate exclusively via HTTP with API-key authentication, with no message queue or event bus mediating inter-service calls. Static service URLs are resolved through Azure App Configuration, and all secrets are managed via Azure Key Vault, consumed at service startup through \texttt{DefaultAzureCredential}.
% TODO: Present the problem definition.


%%------------------------------------------------------------
%% SECTION: Components and Design Decisions
%%------------------------------------------------------------
\section{Proposed Approach}
\label{sec:proposed-approach}

% TODO: Describe the proposed approach.